\clearpage

\section{典型故障与复验}\label{sec:C}

\setcounter{table}{0}

以下记录用于说明车轮方向、入口对准、通道边界、信息补采和对抗调试怎样影响设计选型。整场检验条件和结果见第~\ref{sec:4} 章与第~\ref{sec:adv} 章。

\subsection{姿态符号与轮轴映射}\label{sec:C.1}

两处符号错误曾互相掩护：左右轮接反使转弯方向相反，随后把惯性单元读数取负又使环路成为正反馈。判定方法是让单车以 \ensuremath{\pm}0.8~rad/s 纯自转各跑一回合，裁判由 getOrientation() 的旋转矩阵取 atan2(R10, R00) 作真值航向，控制器同时上报惯性单元与四个编码器的原始值。

\begingroup\small
\setlength{\tableavailable}{\dimexpr\linewidth-8\tabcolsep\relax}
\begin{longtable}{@{}>{\raggedright\arraybackslash}p{0.18000\tableavailable}>{\raggedright\arraybackslash}p{0.17000\tableavailable}>{\raggedright\arraybackslash}p{0.17000\tableavailable}>{\raggedright\arraybackslash}p{0.14000\tableavailable}>{\raggedright\arraybackslash}p{0.34000\tableavailable}@{}}
\caption{纯自转对照的判定性测量}\label{tab:C.1}\\
\toprule
\textbf{指令角速度} & \textbf{裁判真值} & \textbf{惯性单元读数} & \textbf{位置漂移} & \textbf{结论} \\
\midrule
\endfirsthead
\multicolumn{5}{@{}l}{\small\color{muted}表~\thetable\quad （续）}\\
\toprule
\textbf{指令角速度} & \textbf{裁判真值} & \textbf{惯性单元读数} & \textbf{位置漂移} & \textbf{结论} \\
\midrule
\endhead
\midrule
\multicolumn{5}{r@{}}{\footnotesize\color{muted}续下页}\\
\endfoot
\bottomrule
\endlastfoot
+0.8~rad/s & +0.558~rad/s & +0.535~rad/s & 0.001~m & 两者同号，惯性单元为逆时针正 \\
\ensuremath{-}0.8~rad/s & 对称取负 & 对称取负 & 同量级 & 恢复惯性单元读数符号，按横向坐标校正轮轴映射 \\
\end{longtable}
\endgroup


修正后同一单车路线的通道占用由约 23.0~s 降到 7.1~s。相关回合为 \pathref{evidence/webots/} 下的 \pathref{p-att-spin-plus} 与 \pathref{p-att-spin-minus}，单车完整路线闭合见 \pathref{p-motor-b1-route}（D\ensuremath{\to}W\ensuremath{\to}E\ensuremath{\to}L\ensuremath{\to}O，1/1 按期、0 违例）。

\subsection{入口对准的几何检查与回退}\label{sec:C.2}

在这组 Gazebo 场景中，固定对准点到任一队友在途参考线段的距离至少为 0.65~m 时，才采用该对准点。计算方法为对准点到线段的最短距离。历史记录中的一组间距约 0.497~m，其余不小于 0.906~m，阈值用于区分这两类情况。

检查通过时，车辆经固定对准点进入。检查不通过时，车辆改走两点式入口，再核对实际占用、目的地和导航执行条件。0.65~m 适用于本场地与所用参考线段。回退对应记录为 \pathref{physical_entry_alignment_skipped}。

\subsection{口部几何与占用判据的错位}\label{sec:C.3}

早期两次八实体诊断（\pathref{p-motor-full-01} 和 \pathref{p-motor-full-02}）分别完成 11 项和 13 项任务，后一次记录了 5 次占用重叠。原占用窗口为 $2.05\leq x\leq 4.95$~m，走廊两侧墙体沿 $x$ 轴从 1.8~m 延伸至 5.4~m。参考航点 $W(2,0)$ 与 $E(5,0)$ 分别在原窗口西、东边界外约 0.05~m。车体跨过边界时，姿态变化可能使瞬时判定反复切换。

针对这一现象，曾考虑三种办法：将口部余量增到 0.45~m，在 W、E 处移动原航点或增加过渡点，对占用状态切换设置迟滞。这些办法会改变旧回合的判据，故停留在该轮诊断建议。正式整场采用车体包络判空和串行放行，并把对准点移到口部之外，正常回合得到 14/14 按期、5 次进入、5 次清空和 0 次重叠。原始诊断见 \pathref{docs/webots/PHYSICAL_LAYER.md}。

\subsection{问题、修改与复验结果}\label{sec:C.4}

\begingroup\small
\setlength{\tableavailable}{\dimexpr\linewidth-8\tabcolsep\relax}
\begin{longtable}{@{}>{\raggedright\arraybackslash}p{0.12000\tableavailable}>{\raggedright\arraybackslash}p{0.22000\tableavailable}>{\raggedright\arraybackslash}p{0.22000\tableavailable}>{\raggedright\arraybackslash}p{0.22000\tableavailable}>{\raggedright\arraybackslash}p{0.22000\tableavailable}@{}}
\caption{现象、定位、修改与复验}\label{tab:C.2}\\
\toprule
\textbf{层次} & \textbf{现象} & \textbf{定位依据} & \textbf{修改} & \textbf{复验结果} \\
\midrule
\endfirsthead
\multicolumn{5}{@{}l}{\small\color{muted}表~\thetable\quad （续）}\\
\toprule
\textbf{层次} & \textbf{现象} & \textbf{定位依据} & \textbf{修改} & \textbf{复验结果} \\
\midrule
\endhead
\midrule
\multicolumn{5}{r@{}}{\footnotesize\color{muted}续下页}\\
\endfoot
\bottomrule
\endlastfoot
执行器 & 原地打转、航向误差饱和 & \ensuremath{\pm}0.8~rad/s 真值与惯性单元同号 & 去掉航向负号、重排轮轴 & 单车通道占用 23.0~s \ensuremath{\to} 7.1~s \\
资源 & 口部反复占用、5 次重叠 & 占用窗口与航点边界错位 & 车体包络判空、对准外移、驻留迟滞与脱困 & 正常回合：14/14 按期，5 进 5 清，0 重叠 \\
信息 & 过期帧被拒收后 13/14 & 400~s 窗口内无新帧 & 加入主动补采 & 14/14，B2 于 314.240~s 交付 \\
规则 & 后车等到固定 32~s 才放行 & 两车同场景放行时刻对照 & 按前车实际清空放行 & 少等 24.112~s 且无重叠 \\
通信 & 红方无人机约 40~s 停止响应 & 16~ms 高频读观察文件与写入竞争 & 总线读取加共享锁，轮询改为 250~ms & 欺骗回合跑满 400~s，正常完成重放与拦截 \\
对抗 & 虚占至 120~s 时无防卫仍为 14/14 & R1 释放后余量尚足，负对照未显露损失 & 将虚占滞留延长至 145~s & 无防卫降至 12/14，阈值防卫三轮恢复 14/14 \\
\end{longtable}
\endgroup


表~\ref{tab:C.2} 只汇总同平台内的配对结果；跨平台只比较行为是否复现，不并列秒数。
